\documentclass[main.tex]{subfiles}
\begin{document}

\section{Exercise 5 – Two-Stage CMOS Operational Amplifier}

\begin{figure}[H]
    \centering
    \includegraphics[width=0.85\linewidth]{images/exercise5_circuit.png}
    \caption{Two-stage CMOS operational amplifier}
    \label{fig:ex5_circuit}
\end{figure}

\subsection*{Part a) Subcircuit Functionality}

\textbf{Differential Input Pair (M1, M2):} Converts differential input voltage $v_{id} = v_{in1} - v_{in2}$ to differential current $g_{m1}v_{id}$. MOS gates provide essentially infinite input impedance at DC.

\textbf{Active Load (M3, M4, M5):} PMOS current mirror (M3-M4) converts differential current to single-ended output at M2/M4 drain. Provides high gain via high output resistance $r_{ds2} \parallel r_{ds4}$. M5 mirrors reference current to establish input stage bias.

\textbf{Bias Circuit (RBIAS, VIN):} RBIAS sets reference current mirrored throughout circuit. VIN provides gate bias for current sources (M5, M7).

\textbf{Second Stage (M6, M7):} Common-source amplifier with gain $A_{v2} = -g_{m6}(r_{ds6} \parallel r_{ds7})$. M7 acts as current source load.

\textbf{Compensation Network (CC, CB, CA, CL):} Miller capacitor $C_C$ creates pole splitting - dominant pole $p_1$ moves lower, non-dominant pole $p_2$ moves higher, ensuring stability. $C_L$ is load capacitance.

\subsection*{Part b) Parameter Effects on Performance}

\textbf{Increasing Tail Current ($I_{tail}$):}

\textit{Effects:}
\begin{itemize}
    \item Gain: Decreases (intrinsic gain $g_m r_o \propto 1/\sqrt{I_D}$ since $r_o$ drops faster than $g_m$ rises)
    \item GBW: Increases ($GBW = g_{m1}/(2\pi C_C)$ and $g_m \propto \sqrt{I_D}$)
    \item Slew Rate: Increases ($SR = I_{tail}/C_C$)
    \item Stability: Improves (non-dominant pole $p_2 \approx g_{m6}/C_L$ moves higher)
\end{itemize}

\textit{Pros:} Higher speed, better slew rate, improved phase margin

\textit{Cons:} Higher power consumption, lower gain, increased noise

\textbf{Increasing W/L of Input Pair (M1, M2):}

For fixed $I_D$: increases $g_{m1}$ (∝ $\sqrt{W/L}$), thus increasing first-stage gain and GBW. Reduces input-referred noise (∝ $1/g_m$). Improves matching.

\textit{Cons:} Larger parasitic capacitances, increased area

\textbf{Increasing W/L of Second Stage (M6):}

Increases $g_{m6}$, pushing $p_2 = g_{m6}/C_L$ higher, improving phase margin. Enables higher output current drive and better slew rate.

\textit{Cons:} Increased area and parasitics

\textbf{Increasing Compensation Capacitor ($C_C$):}

\textit{Stability:} Lowers dominant pole $\omega_{p1} \approx 1/(R_{out,1} \cdot A_{v2} \cdot C_C)$, improving phase margin by separating poles further

\textit{GBW:} Decreases ($GBW = g_{m1}/(2\pi C_C)$)

\textit{Pros:} Better stability, reduced overshoot, higher phase margin

\textit{Cons:} Reduced bandwidth, slower settling, larger area

\subsection*{Design Trade-offs}

Fundamental trade-offs: Power vs. Speed (higher current = faster but more power), Gain vs. Bandwidth (fixed GBW product), Stability vs. Performance (larger $C_C$ = stable but slow). Optimal design balances these based on application requirements.

\end{document}
